\documentclass[11pt,a4paper]{article}
\usepackage[margin=23mm]{geometry}
\usepackage{amsmath,amssymb,amsthm,hyperref}
\newtheorem{lemma}{Lemma}
\newcommand{\gen}[1]{\langle#1\rangle}
\title{The two-row endpoint of the simultaneous-swap argument}
\author{Human Compression project: bounded L044 proof component}
\date{}
\begin{document}\maketitle
We use the complete accepted block-deficit and small-endpoint
proofs, including their L047/Whiston hypotheses and dependencies.
Thus $G\leq A_n$ is a transitive string C-group with every
maximal omission intransitive, a unique eligible $i$-edge joins
the two $G_i$-orbits, and each of those orbits has two blocks.
Put $s=\rho_i$. Products below act from left to right.

\begin{lemma}
If one block size is two, $n\geq12$, and $\alpha=1$, then
$r\leq\lfloor(n-1)/2\rfloor$. Consequently, together with the
accepted small-endpoint proof, the simultaneous-swap case
$m_1=m_2=2$, $\alpha=1$ is completely excluded as an obstruction
to this rank bound.
\end{lemma}
\begin{proof}
Write the two block sizes as $2,b$, where $b\geq4$.
A failure gives $\beta=\gamma=0$, $r=b+2$, and a single
simultaneous fixed-point-free swap $\tau=\rho_{i-1}$ after
reversal. Its evenness forces $b$ even. The complete block
image of $G_i$ is the diagonal $C_2$ generated by $\tau$.

Let $x=\rho_{i-2}$ if it exists, and let $C$ be all other
distinguished generators of $G_i$ apart from $\tau,x$.
They centralize $\tau$. As in the accepted proof, row projection
of $\gen C$ is faithful: its kernel is contained in $\gen\tau$
and IP makes this intersection trivial. The group $\gen C$
is a character graph over a subgroup of $S_2\times S_b$.
Except for $c=\rho_{i+1}$, its generators fix both row anchors
and have trivial block character. The independent-length bound
in $S_1\times S_{b-1}$ is $b-2$. It forces both $x$ and $c$
to exist, gives $|C|=b-1$, and identifies
\[
 H_0=\gen{C\setminus\{c\}}=S_{b-1}
\]
diagonally on the large orbit's nonanchor rows, fixing all four
points of the small orbit.

The generators of $H_0$ split into the prefix at indices at
most $i-3$ and the tail at indices at least $i+2$.
These subgroups commute and have trivial intersection by IP.
Since the natural $S_{b-1}$ is centerless and directly
indecomposable for $b-1\geq3$, exactly one interval is nonempty.
If the prefix is empty, $x$ centralizes all of $H_0$.
It preserves each large block: it fixes its crossing endpoint
and its block image is either the identity on both systems or
the simultaneous exchange. The natural $S_{b-1}$ has one fixed
point in each $b$-point block and trivial centralizer there,
so $x$ fixes the entire large orbit. On the small orbit it
fixes the two-point block containing its crossing endpoint,
and can perform at most one transposition on the other block.
It is therefore the identity or odd, impossible.

Hence all of $H_0$ is the left prefix, and $c$ is the sole
right-interval generator. It centralizes $H_0$, so its large
row projection is trivial. Faithfulness of row projection
forces its small row projection to be the transposition;
otherwise $c$ would be the identity. Thus the row group is
$S_2\times S_{b-1}$, and the character of $c$ is the only
remaining possible nontrivial block character.
If that character is nontrivial, $c$ exchanges the two large
blocks row by row. Since $x$ commutes with $c$ and fixes the
large crossing endpoint, it also fixes its opposite row point.
The prefix fixes both, and $\tau,c$ only exchange them.
Their pair is then a $G_i$-invariant proper subset of the large
orbit, a contradiction. Hence this character is trivial.

Label the small orbit $A_0,A_1,A'_0,A'_1$ and the large
blocks $0,1,\ldots,b-1$ and $0',1',\ldots,(b-1)'$,
with crossing endpoints $A_0,0$. We have
\[
 c=(A_0\ A_1)(A'_0\ A'_1),\qquad
 \tau=(A_0\ A'_0)(A_1\ A'_1)
                  \prod_{j=0}^{b-1}(j\ j').                \tag{1}
\]
The prefix is a maximal independent string of $S_{b-1}$,
and $x$ centralizes its final-end omission.

First take even $b\geq6$. The accepted L047 classification
makes this prefix an ordinary Coxeter string. Its exceptional
outer degree-six action does not occur, since $b-1$ is odd.
After relabeling, let its final generator be
\[
 p=(b-2\ b-1)((b-2)'\ (b-1)').
\]
The final-end omission is $S_{b-2}$ on rows $1,\ldots,b-2$,
fixing $0,b-1$. Thus $x$ is the identity on the first large
block, since it fixes $0$, and its second-block action is
either the identity or $(0'\ (b-1)')$. The latter is necessary
to connect the large anchor pair to the other rows. On the
small orbit $x$ can only swap $A'_0,A'_1$. Evenness now forces
\[
                   x=(A'_0\ A'_1)(0'\ (b-1)').             \tag{2}
\]
These literal permutations satisfy
\[
                  c=\bigl[\tau(xp)^2x\bigr]^2p.           \tag{3}
\]
Indeed, putting $u=b-2$ and $v=b-1$, the bracketed word has
cycles
\[
 (A_0\ A'_1\ A_1\ A'_0),\qquad
 (0\ 0'),\qquad(u\ v'\ v\ u')
\]
on these ten points. Its square is $cp$. On every other row
pair it is just $\tau$, whose square is the identity.
Equation (3) makes the distinguished generator $c$ redundant.

It remains to treat $b=4$ explicitly. The prefix is
\[
 p=(1\ 2)(1'\ 2'),\qquad q=(2\ 3)(2'\ 3'),
\]
in that order before $x,\tau,s,c$. The centralizer of $p$
on each four-point block gives
\[
 x=(A'_0\ A'_1)^\eta(1\ 2)^\epsilon
                (1'\ 2')^\delta(0'\ 3'),\qquad
 \eta+\epsilon+\delta=1\pmod2.                            \tag{4}
\]
The factor $(0'\ 3')$ is necessary for the large orbit's
transitivity, and fixation of the two crossing endpoints
excludes every other factor. If $\epsilon=\delta$, then
$\eta=1$ and $x_0=xp^\epsilon=(A'_0\ A'_1)(0'\ 3')$.
The same ten-point calculation as (3), now using the final
prefix generator $q$, gives
\[
                  c=\bigl[\tau(x_0q)^2x_0\bigr]^2q.       \tag{5}
\]
The word uses only distinguished generators other than $c$,
so both of these cases are excluded.

If $\epsilon\ne\delta$, then $\eta=0$ and $x$ fixes the
small orbit pointwise. The subgroup $\gen{p,q,x}$ has on the
large orbit the three point orbits $\{0\}$, $\{1,2,3\}$ and
$\{0',1',2',3'\}$, with natural symmetric projections on the
last two. As $s$ centralizes this subgroup, it fixes those
seven points individually. Its crossing action and evenness
therefore give
\[
 s=(A_0\ 0)h,\qquad
 h\in\{(A_1\ A'_0),(A_1\ A'_1),(A'_0\ A'_1)\}.           \tag{6}
\]
Put $w=s(\tau s x)^5\tau x$. If $h=(A_1\ A'_1)$, the
ten-point identity of the accepted four-row endpoint gives
$w=p$. For the other two choices one has
\[
                         p=(\tau s\tau w)^3.              \tag{7}
\]
For a direct check, $\tau s\tau w$ is the product of $p$
and the following two disjoint three-cycles:
\[
\begin{array}{c|l}
h&\text{three-cycles}\\\hline
(A_1\ A'_0)&(A_0\ 0\ 3),\quad(A_1\ A'_0\ 3'),\\
(A'_0\ A'_1)&(A_0\ 0\ 3),\quad(A'_0\ 3'\ A'_1).
\end{array}
\]
These formulas hold for both choices $(\epsilon,\delta)$.
Cubing gives (7). Thus all six remaining choices make $p$
redundant. This completes the case $b=4$ and the proof.
\end{proof}

\paragraph{Scope and attribution.}
This is a further bounded proof of the both-imprimitive rank
branch in L044, Section~5.1, physical pages~23--25. It retains
the exact common-theory and source attribution of the accepted
block-deficit note. The explicit word identities above are
new proofs for the residual endpoint, with complete cycle
verification; bounded word-finding experiments are not
required premises. Together with the accepted endpoint note,
they complete the $\alpha=1$ two-block case. They do not
settle $\alpha=0$, mixed block-count configurations, or all of
L044. There is no new original or conjecture credit and no
new formal declaration.
\end{document}
